\paragraph{What the model saw.} Each revision request reproduced the original specification and included only the active sizes, their incumbent weights, and diagnostics. For each size and split the model received the scaled loss, raw RMSE, whole-example exact count, and at most three non-exact inputs with the largest absolute scaled error on any coordinate. For each of those inputs it saw the target and current prediction. The 256 central development inputs were fixed within task and start and shared across sizes and prompt conditions. Integer indices 9--16 instead extended the original nine labels. Original one-shot replies, new prompts, replies, decisions and final vectors are retained. Appendix~\ref{app:repair} is a different experiment that allows a circuit description to change.
\paragraph{Stopping and selection.} The collector made \revCalls{} revision requests after the original 60 model calls. Their recorded median latency was \latencySec{} seconds per request, with \outputTokens{} output tokens in total; requests for different tasks ran concurrently. These requests produced 332 size-specific proposals, of which 217 were accepted; one width-16 square proposal had the wrong number of coefficients. None of the 161 requests used tools. The 174 size-by-condition trajectories share requests within an original task, start and prompt condition. Of those trajectories, \zeroStops{} were exact before any revision. At the end, \exactStops{} had stopped on exact training and development outputs, \patienceStops{} after two consecutive rejected proposals, and \budgetStops{} at the three-attempt cap. After all selections were frozen, a separate evaluator reconstructed the feedback from training and development data, checked every decision against the saved replies and then scored the selected weights. Table~\ref{tab:iterrounds} follows all 87 settings per condition through each revision opportunity. The cap limits what we can conclude about additional rounds.
\paragraph{What if gradient descent can use all feedback labels?} We separately fit both gradient procedures from the same Xavier start using the original training observations pooled with every feedback label, keeping both optimizer budgets and the test sets fixed. This gives the gradient solver 268 labeled examples per sampled task or 17 integer indices; Astra saw diagnostics and selected examples from that pool. The baseline therefore has access to more individual labels than Astra. With examples only, iterative WSP had lower outer RMSE than pooled-data gradient descent in \vsPoolGD{} settings and than its L-BFGS refinement in \vsPoolGL{} (out of 87). The pooled procedures matched every outer reference in \perfectPoolGD{} and \perfectPoolGL{} settings, respectively. Machine-readable files include training, feedback, independently held-out central or early, and outer scores for both pooled endpoints.
\paragraph{Independent test draw.} After the weight selection was frozen, we generated a second set of central and outer inputs using a new random seed. On the 69 settings with sampled inputs, the iterative-versus-one-shot outer RMSE comparison was 22 lower, 23 tied and 24 higher with examples only, and 20 lower, 23 tied and 26 higher with the formula supplied. The 18 integer settings have fixed target indices and are not repeated as independent draws. The companion data include the scores from both draws.

\begin{table}[htbp]\centering\small
\caption{What changed with each opportunity to revise? Each information condition keeps all 87 settings throughout. Attempted and accepted count individual size-specific proposals. After a size stops, its last retained weights carry forward. The last four columns are retrospective extrapolation scores; these scores were never shown to Astra. Exact counts settings in which every test prediction matched. Lower, equal and higher compare the retained outer RMSE with that setting\textquotesingle s one-shot outer RMSE.}\label{tab:iterrounds}
\begin{tabular}{@{}lrrrrrrr@{}}\toprule
Information & Round & Attempted & Accepted & Exact & Lower & Equal & Higher\\\midrule
Examples only & 0 & -- & -- & 30 & 0 & 87 & 0\\
 & 1 & 57 & 51 & 30 & 24 & 36 & 27\\
 & 2 & 57 & 32 & 30 & 28 & 35 & 24\\
 & 3 & 52 & 22 & 30 & 28 & 35 & 24\\
\midrule
Formula and examples & 0 & -- & -- & 30 & 0 & 87 & 0\\
 & 1 & 57 & 51 & 30 & 27 & 36 & 24\\
 & 2 & 57 & 34 & 30 & 27 & 35 & 25\\
 & 3 & 52 & 27 & 30 & 27 & 35 & 25\\
\bottomrule\end{tabular}\end{table}


\input{diagnosis}
\begin{figure}[htbp]\centering
\includegraphics[width=\linewidth]{square_continuation.pdf}
\caption{The actual width-16 square networks from the examples-only condition at random start 164. Training and development inputs lie in the central interval. The stored one-shot weights retain changes of slope outside that interval. Iteration concentrates all changes of slope inside it and extends an affine function into each outer interval. The curves explain this specific extrapolation regression.}\label{fig:square-continuation}
\end{figure}

